\scanpage{10}\mspage{8}
\[
\mathrm{Hom}_{M_{0\pi}}(a,b)\simeq\pi_0\bigl(\mathrm{Ch}(a,b)\bigr).
\]
\deleted{\unsure{Cf. cette}} Cette relation est la \uline{définition} de $\mathrm{Hom}_{M_{0\pi}}$, mais ici la relation est entendue comme une bijection \uline{fonctorielle} en $a,b$.

On trouve notamment que si
\[
u:a'\unsure{\text{---}|\longrightarrow}a,\qquad v:b\ \text{---}\ b'
\]
sont des flèches \uline{inversibles} \inserted{\unsure{de|dans} $M_{0\pi}$}, alors
\[
\mathrm{Ch}_{\deleted{\unsure{\infty}}}(a,b)\longrightarrow\mathrm{Ch}(a',b')
\]
correspondant dans $\unsure{\mathscr H|\mathcal H}$ est \uline{iso}.

On appliquera tout ceci \unsure{surtout} \unsure{aux} deux situations déduites d'une situation $(M,\Sigma)$, avec
\[
\begin{aligned}
1^{\mathrm o})\quad &\mathscr H=\mathrm{Hot},&
 \mathrm{Ch}(a,b)&=\mathrm{Ch}^{\Sigma}_{\infty\deleted{\unsure{\Sigma}}}(M;a,b)\\
2^{\mathrm o})\quad &\mathscr H=\inserted{\mathrm{Ind}}\,\mathrm{Cat}\,\mathrm{hot},&
 \mathrm{Ch}(a,b)&=\mathrm{CH}^{\Sigma}_{\infty}(M;a,b).
\end{aligned}
\]
% slanted
\margin{$*$ dans $\mathrm{Hot}$ et $\mathrm{Ind}\,\mathrm{Cat}\,\mathrm{hot}$ est le produit cartésien sans plus}
Dans les deux cas, on a donc \deleted{\unsure{un}} un foncteur sur $(M\Sigma^{-1})^{o}\times(M\Sigma^{-1})$, à valeurs soit dans $\mathrm{Hot}$, soit dans \deleted{$\unsure{\mathrm{Cat}}$} $\mathrm{Ind}\,\mathrm{Cat}\,\mathrm{hot}$, (\uline{NB} les deux structures de cat. sur $M_0$ déduites de ces deux foncteurs sont can\textsuperscript{t} isomorphes)
\diagram{\begin{tikzcd}[ampersand replacement=\&,row sep=small,column sep=small]
\& \mathrm{Hot} \& (a,b)\unsure{\longrightarrow|\twoheadrightarrow}\mathrm{Ch}_{\infty}^{\Sigma}(a,b)\\
(M\Sigma^{-1})^o\times(M\Sigma^{-1}) \arrow[ur] \arrow[dr] \& \& \\
\& \mathrm{Ind}\,\mathrm{Cat}\,\mathrm{hot} \& (a,b)\ \text{---}\ \mathrm{CH}_{\infty}^{\Sigma}(a,b)\unsure{,|.}
\end{tikzcd}}
% archive sheet number: 9

